\documentclass[12pt,a4paper]{article}

\usepackage[T1]{fontenc}
\usepackage[utf8]{inputenc}
\usepackage[polish]{babel}
\usepackage{lmodern}
\usepackage{microtype}
\usepackage{geometry}

\geometry{margin=2.5cm}
\setlength{\parindent}{0pt}
\setlength{\parskip}{0.7em}
\pagestyle{empty}

\begin{document}

\section*{Streszczenie}

Praca dotyczy dostrajania małych modeli językowych z otwartymi wagami do
binarnej klasyfikacji wiadomości jako pożądane albo niepożądane. Po przeglądzie
dostępnych modeli do badań wybrano Qwen3-0.6B. Wykorzystano metodę niskorangowej
adaptacji LoRA. Porównano 16 konfiguracji oraz stany modelu
zapisane w kolejnych etapach treningu. Następnie sprawdzono cztery sposoby
użycia modelu jako klasyfikatora. Porównanie objęło także model bez dostrajania
i sześć innych klasyfikatorów. Końcową ocenę przeprowadzono na czterech
rozłącznych zbiorach testowych, obejmujących po 2000 wiadomości. Na czterech
zbiorach testowych najlepszy dostrojony wariant osiągnął średnią dokładność
klasyfikacji równą 97,8\%. DistilBERT, najskuteczniejszy z pozostałych
klasyfikatorów, osiągnął średnią dokładność 94,4\% i działał około 2,3 raza
szybciej od Qwen3.
Zastosowanie dostrojonego modelu w filtrze poczty wymaga uwzględnienia kosztu
obliczeniowego klasyfikacji i sprawdzenia go na bieżącym ruchu pocztowym.

\textbf{Słowa kluczowe:} niechciana poczta elektroniczna, klasyfikacja spamu,
modele językowe, Qwen3, LoRA

\end{document}
